\section{The standalone indexer}\label{sec:indexer}

\subsection{Crate boundaries}

The Rust implementation is a three-crate workspace.

\begin{table}[htbp]
\centering
\small
\caption{Indexer crates and responsibilities.}
\label{tab:crates}
\begin{tabularx}{\linewidth}{@{}p{0.24\linewidth}X@{}}
\toprule
\textbf{Crate} & \textbf{Responsibility} \\
\midrule
\code{hopper-core} & Pure protocol: envelope, v8 and legacy body, scripts, parameters,
generations, wrap, sat math, art, PNG, fold document. Depends only on
\code{sha2} (plus optional serde). \\
\code{hopper-indexer} & Pure state machine: sat tracker, rules, timers,
events, snapshots, optional redb store. Performs no network I/O; the optional
redb store is local disk persistence. \\
\code{hopper-indexer-api} & Axum HTTP API and block ingest (ZMQ/RPC or admin
POST). Owns the network and the write lock. \\
\bottomrule
\end{tabularx}
\end{table}

The separation is not cosmetic. \code{hopper-core} can be compiled for
native and mobile targets without a Bitcoin node. \code{hopper-indexer} can
be replayed from a fixture without sockets. Only the API crate knows about
ZMQ, RPC, HTTP, or authentication.

\subsection{Ingest}

The indexer always ingests from block 0. A later start cannot name coins
created earlier; a frog whose mint predates the start would look as if it had
been sent away. This is a correctness rule, not an optimization.

There are two wires, sharing the same block decoder and the same write lock.

\begin{itemize}[leftmargin=1.4em,itemsep=0.3em]
  \item \code{POST /v1/admin/block} accepts an operator or debug JSON block.
        Outside loopback it requires the admin token.
  \item bitcoind ZMQ \code{hashblock} plus RPC is the production path. On a
        new hash, the indexer fetches blocks from its tip to that hash. On a
        reorg, it rolls back to the fork and replays.
\end{itemize}

The ingest layer converts a Bitcoin transaction into a small, witness-light
structure:
\begin{lstlisting}[language={},caption={Ingest types.},label={lst:ingest}]
Tx {
    txid,
    inputs[{ prevout, value, inscription? }],
    outputs[{ value, op_return? }]
}
Block { hash, height, txs }
\end{lstlisting}
It parses the tapscript witness at \code{witness[n-2]} with the ord envelope
decoder, extracts the control block at \code{witness[n-1]}, and records the
content type, body and tag-3 parents. The state machine never touches
witnesses.

\begin{figure}[htbp]
\centering
\begin{tikzpicture}[
  box/.style={draw=deepgreen, fill=pale, rectangle, align=center,
              text width=25mm, minimum height=13mm, inner sep=3pt,
              font=\scriptsize\sffamily},
  net/.style={draw=midgreen, fill=white, rectangle, align=center,
              text width=25mm, minimum height=13mm, inner sep=3pt,
              font=\scriptsize\sffamily},
  arrow/.style={-{Stealth[length=2mm]}, draw=deepgreen!75, line width=0.7pt},
  lbl/.style={font=\tiny\sffamily, color=muted, fill=white, inner sep=1pt}
]
\node[net] (node) at (0,0) {bitcoind\\ZMQ + RPC};
\node[box] (ingest) at (3.6,0) {ingest\\decode + verify};
\node[box] (engine) at (7.2,0) {state machine\\frogs, coins, wrap};
\node[box] (store) at (10.8,0) {snapshot\\events / redb};
\node[box] (api) at (7.2,-2.2) {HTTP API\\axum};
\node[net] (clients) at (7.2,-4.0) {wallet / PWA / ord};
\draw[arrow] (node) -- node[lbl, above] {blocks} (ingest);
\draw[arrow] (ingest) -- node[lbl, above] {Tx/Block} (engine);
\draw[arrow] (engine) -- node[lbl, above] {state} (store);
\draw[arrow] (engine) -- (api);
\draw[arrow] (api) -- (clients);
\end{tikzpicture}
\caption{Indexer architecture. The state machine is independent of the
network and of the HTTP layer; persistence is a projection of the same
state.}
\label{fig:indexer}
\end{figure}

\subsection{State, snapshots and the state root}

The state is a function
\begin{equation}
(\mathrm{config},\mathrm{seeds},\mathrm{chain})\;\longrightarrow\;
\mathrm{IndexState}+\mathrm{event\ log}.
\end{equation}
It contains the frog records, the coin ranges, the retained headers, the
pending timers and the tip. A canonical state root hashes the parameter
commitment, the tip (height and hash), the ordered frog records and the coin
ranges. Two indexers that agree on consensus must agree on the root; a
disagreement is a bug or a different parameter set, not a client preference.

Persistence has three layers.

\begin{itemize}[leftmargin=1.4em,itemsep=0.3em]
  \item A binary snapshot at \code{HOPPER_SNAPSHOT_PATH}, default
        \code{hopper-indexer.snapshot.bin}. The snapshot embeds the parameter
        hash; a mismatch on load is a hard error, never a migration.
  \item Append-only event segments next to the snapshot.
  \item An optional redb durable store (feature \code{redb}, path
        \code{HOPPER_DB}) holding the full state and the undo journals.
\end{itemize}

The snapshot interval defaults to 10 blocks, the retained block window to
288 blocks, and the event segment size to 4096 events.

\subsection{HTTP API}

The public surface is read-only. Administrative routes are separate and
require a Bearer token when \code{HOPPER_ADMIN_TOKEN} is set; the server
refuses a non-loopback bind without one. They are never exposed by the PWA.

\begin{table}[H]
\centering
\small
\caption{Indexer HTTP API.}
\label{tab:api}
\begin{tabularx}{\linewidth}{@{}p{0.39\linewidth}X@{}}
\toprule
\textbf{Route} & \textbf{Purpose} \\
\midrule
\code{GET /v1/status} & Tip, height, active crossings, wire version \\
\code{GET /v1/params} & Full parameter block, \code{params_hash}, chain id \\
\code{GET /v1/generation/{height}} & Generation index, start, end,
\code{required(h)} \\
\code{GET /v1/wrap/{height}} & The block's wrap ($0|1|2$); 404 until
ingested \\
\code{GET /v1/froggers} & Paginated list; \code{limit}, \code{offset},
\code{after_sat}, \code{to_sat} \\
\code{GET /v1/frogger/{sat}} & One frog record, with \code{gen},
\code{in_flight.dir} and \code{outcome.wrap} \\
\code{GET /v1/road} & In-flight crossings: $T$, direction, target, wrap,
settle height \\
\code{GET /v1/leaderboard} & Ranked by generation, then \code{dist};
\code{limit} \\
\code{GET /v1/state/root} & Canonical state root \\
\code{GET /v1/events} & Event log; \code{limit} \\
\code{GET /v1/events/stream} & Server-sent event stream \\
\code{GET /v1/mempool} & Mempool preview; needs RPC \\
\code{GET /v1/fees} & bitcoind fee estimate; never invented \\
\code{GET /v1/utxo/{txid}/{vout}} & Value, ordinal ranges, frogs,
\code{exact} flag \\
\code{GET /healthz}, \code{/readyz}, \code{/metrics} & Liveness, readiness,
metrics \\
\code{POST /v1/admin/block} & Ingest an operator/debug block \\
\code{POST /v1/admin/seed} & Bootstrap explicit ordinal ranges \\
\code{POST /v1/admin/rollback} & Roll back to a fork point \\
\bottomrule
\end{tabularx}
\end{table}

CORS is a browser read policy configured by \code{HOPPER_CORS_ORIGINS}. It
is not access control and must not be treated as one. The API never publishes
a wrap for a height it has not ingested: the outcome of a future block is not
predictable from the indexer either.

\subsection{Configuration}

\begin{table}[H]
\centering
\small
\caption{Indexer configuration. Legacy \code{FROGGER_*} names are accepted
as fallbacks.}
\label{tab:config}
\begin{tabularx}{\linewidth}{@{}p{0.43\linewidth}X@{}}
\toprule
\textbf{Variable} & \textbf{Meaning / default} \\
\midrule
\code{HOPPER_API_ADDR} & Bind address; default \code{127.0.0.1:9770} \\
\code{HOPPER_ACTIVATION_HEIGHT} & Frozen inscription cutover; default and
required \code{974,034} \\
\code{HOPPER_INSCRIPTION_ACTIVATION_HEIGHT} & Frozen cutover; default and
required \code{974,034} \\
\code{HOPPER_RATE_LIMIT_PER_MIN} & Public-read budget per caller; default
300, \code{0} disables \\
\code{HOPPER_CORS_ORIGINS} & Browser read origins; default \code{*}; not
access control \\
\code{HOPPER_SNAPSHOT_PATH} & Snapshot file; default
\code{hopper-indexer.snapshot.bin} \\
\code{HOPPER_SNAPSHOT_EVERY} & Snapshot interval; default 10 blocks \\
\code{HOPPER_RETAIN_BLOCKS} & Reorg window; default 288 \\
\code{HOPPER_EVENT_SEGMENT} & Events per append-only segment; default 4096 \\
\code{HOPPER_EVENT_RETAIN_SEGMENTS} & Old segments retained; default 0 (keep
all) \\
\code{HOPPER_SEEDS_FILE} & Bootstrap seed JSON for fresh starts \\
\code{HOPPER_INGEST_START_HEIGHT} & Must be 0 or unset; later starts are
refused \\
\code{HOPPER_ORD_URL} & Optional ord base for lazy range seeding \\
\code{HOPPER_ZMQ_ADDR} & bitcoind \code{zmqpubhashblock}, e.g.
\code{tcp://127.0.0.1:28332} \\
\code{HOPPER_RPC_URL} & bitcoind RPC; default
\code{http://127.0.0.1:18443} \\
\code{HOPPER_RPC_COOKIE}\newline \code{HOPPER_RPC_USER}\newline \code{HOPPER_RPC_PASS} & RPC authentication \\
\code{HOPPER_NETWORK} / \code{HOPPER_GENESIS_HASH} & Published chain
identity \\
\code{HOPPER_ADMIN_TOKEN} & Bearer token for \code{/v1/admin/*}; a
non-loopback bind without it is refused \\
\code{HOPPER_DB} & Optional redb path \\
\code{HOPPER_DB_CACHE_MIB} & redb page cache in MiB; 0/unset uses redb's
default \\
\code{HOPPER_ENGINE=store} & Keep only a block-sized overlay in memory \\
\code{HOPPER_REINDEX=1} & Discard the durable store \\
\bottomrule
\end{tabularx}
\end{table}

The legacy \code{FROGGER_*} spellings are read as fallbacks so an existing
deployment can migrate without editing its service file.

\subsection{Seeding}

Ordinal ranges ultimately come from the chain. A fresh indexer ingests from
genesis, but a deployment that starts from a snapshot, or that is handed a
UTXO it never saw move, can be given explicit ranges:
\begin{center}
\begin{minipage}{0.88\linewidth}
\begin{lstlisting}[language={},caption={Admin seed request.},label={lst:seed}]
POST /v1/admin/seed
{
  "utxos": [
    { "txid": "<display hex>", "vout": 0,
      "ranges": [[start, end]] }
  ]
}
\end{lstlisting}
\end{minipage}
\end{center}
A coin whose ranges were never observed is a shadow sat. A mint on one is
rejected as \texttt{MintAnchor\allowbreak Untracked}. Seeding is an operator bring-up step,
not a consensus rule.

\subsection{Rejections}

No malformed input is chain-invalidating. A rejection leaves state unchanged
while sat movement still applies. The closed set lives in
\code{hopper-indexer/src/types.rs} as \code{RejectionReason}; its families
are:
\begin{itemize}[leftmargin=1.4em,itemsep=0.2em]
  \item mint anchor issues (untracked, shadow, already carrying a frog);
  \item inscription content type and mint-body digest issues;
  \item crossing body, parent, leaf or control-block issues;
  \item the direction set: \code{dir > 2}, the spent leaf is not the
        \code{dir} leaf, or the control block is not the matching row of
        Table~\ref{tab:control};
  \item crossing outpoint and output-position issues;
  \item crossing on a non-idle frog;
  \item funding inputs carrying a frog;
  \item mixed \code{FRG} envelope;
  \item parameter-hash mismatch.
\end{itemize}

A rejected transaction is not a reorg and not a chain failure. It is a
protocol message that the state machine declined to apply. The API also
refuses to serve a mismatched parameter set: \code{GET /v1/params} publishes
the hash, wire version and chain identity, and the wallet refuses to sign
without a matching \code{network}. Missing identity is not permission.

\begin{invariant}{Deterministic replay}
Given the same blocks, the same seeds and the same parameter set, the
indexer produces the same state root and the same event log. Network
topology, arrival order and wall-clock time are not inputs.
\end{invariant}
